% 张三慧《大学物理学》第四版第二卷，第17章《磁场和它的源》
% 原教材扫描 PDF 物理页142--144（印刷125--127页）。
% 17.1--17.20为章末习题，必须与上一页的思考题17.x区别。
% 图17.38--17.46均属于源题物理条件；独立矢量重绘/同行核验仍待完成。
\section*{第17章\quad 磁场和它的源}
\subsection*{章末习题17.1--17.20}

\paragraph{习题 17.1}
求图17.38(a)(b)(c)各图所示
\(P\) 点的磁感应强度 \(B\) 的大小和方向。
(a) 两条互相垂直的半无限直导线在直角处连接，
电流沿竖直段向下，沿水平段向左，
\(P\) 在水平段向右的延长线上，距交点 \(a\)。
(b) 上下两条水平的半无限直导线用右侧
半径为 \(r\) 的半圆弧连接，电流从上方
向右流进、沿右侧半圆顺时针向下、再沿下方
向左流出，\(P\) 在半圆中心。
(c) 边长 \(a\) 的等边三角形闭合导线，
电流沿图示方向顺时针流动，
\(P\) 在正三角形中心。

\paragraph{习题 17.2}
高压输电线在地面上空
\(25\,\mathrm m\) 处通过电流
\(1.8\times10^3\,\mathrm A\)。
（1）求在正下方地面处由输电线产生的
磁感应强度；
（2）在此地区地磁场约为
\(0.6\times10^{-4}\,\mathrm T\)，
输电线产生的磁场与地磁场相比如何？

\paragraph{习题 17.3}
汽船上指南针距载流导线
\(0.80\,\mathrm m\)，
导线电流为 \(20\,\mathrm A\)。
（1）将导线近似为长直导线，
求它在指南针处的磁感应强度；
（2）当地地磁水平分量（指向北）
为 \(0.18\times10^{-4}\,\mathrm T\)，
若导线磁场水平且与地磁场垂直，
求指南针偏离正北的角度，
并求这种偏角所能达到的最大值。
\textit{须说明几何方向对偏角的影响。}

\paragraph{习题 17.4}
两根导线沿半径方向引到铁环上
\(A,C\) 两点（图17.39），
其电流方向如图。
电流由 \(A\) 处分流，
沿环的两条弧分别到达 \(C\) 重新汇合。
求环中心 \(O\) 处的磁感应强度大小。
假设导线、环粗细均匀，同种材料，
忽略径向引线在圆心处的磁贡献。

\paragraph{习题 17.5}
两条互相平行的长直导线间距
\(d=40\,\mathrm{cm}\)，
流过相反方向的等大电流
\(I_1=I_2=20\,\mathrm A\)，见图17.40。
（1）求在两导线所在平面内、
与两导线等距的中线上一点的磁感应强度；
（2）求穿过图中灰色矩形区域的磁通量，
其离左、右线距离
\(r_1=r_3=10\,\mathrm{cm}\)，
沿导线方向的长度
\(l=25\,\mathrm{cm}\)。
\textit{原图左线电流向下、右线向上。}

\paragraph{习题 17.6}
半径 \(R\) 的半圆形载流导线，
平面垂直于半圆轴线，
电流为 \(I\)。
求在轴线上离圆心
\(x\) 处的磁感应强度
\(\boldsymbol B\)（图17.41）。
\textit{半圆导线不是闭合整圆，
按毕奥—萨伐尔积分须考虑磁场的横向分量。}

\paragraph{习题 17.7}
一根电磁铁的大电流长引线是同轴结构：
中央为直径
\(5.0\,\mathrm{cm}\) 的实心铝棒，
外面套内径 \(7.0\,\mathrm{cm}\)、
外径 \(9.0\,\mathrm{cm}\) 的铝圆筒作为电流回路，
中间填充不导电的冷却油。
进、回电流大小都为
\(I=5.0\times10^3\,\mathrm A\)，
方向相反。
设两部分导体横截面各处电流密度均匀，
求距轴线 \(r\) 处磁场分布并画 \(B-r\) 曲线。
\textit{忽略铝与油自身磁化效应。}

\paragraph{习题 17.8}
由教材长直载流导线磁场公式，
用积分法推导：
（1）无限长、横截面均匀通电的实心圆柱
内部和外部磁感应强度；
（2）无限大均匀面电流（面电流密度
\(j\)，单位 \(\mathrm{A/m}\)）
两侧的磁场方向与大小。

\paragraph{习题 17.9}
按图17.11讨论半无限和无限长载流直导线，
沿导线从 \(-\infty\) 到 \(+\infty\)
对毕奥—萨伐尔定律积分，
独立证明无限直线电流的
\(B=\mu_0I/(2\pi r)\)，
与教材式(17.24)比较。

\paragraph{习题 17.10}
设一个足够长的均匀密绕螺线管，
每单位长度匝数为 \(n\)，通电流 \(I\)。
构造图17.42所示的安培矩形回路，
利用安培环路定理推导管内
\[
 B=\mu_0nI.
\]
说明管外近似磁场和端部效应的限制。

\paragraph{习题 17.11}
两段半无限长直螺线管紧密对接，
组成无限长直螺线管；
对图17.43所示\emph{单根半无限长螺线管}
用叠加、磁感线连续性证明：
（1）通过管口截面的磁通量是远离管口时
管内截面的磁通量的一半；
（2）紧靠管口的磁感线
\(abc\) 在管外段为垂直于管轴的直线；
（3）从管侧面漏出的磁感线向管外弯离，
如 \(def\) 线所示；
（4）在管内远离轴线 \(r_0\) 的磁感线，
穿出管口时到轴的距离
\(r=\sqrt2\,r_0\)。
\textit{原图及口部场分布是证明所需条件，
不能只套无限长螺线管内场。}

\paragraph{习题 17.12}
在托卡马克装置中，
环管轴线半径
\(2.0\,\mathrm m\)，
环管横截面半径
\(1.0\,\mathrm m\)；
圆环上均匀缠绕总长
\(10\,\mathrm{km}\) 的水冷铜线圈，
线圈中流过峰值
\(7.3\times10^4\,\mathrm A\)
的脉冲电流。
在近似按安培环路定理计算的模型中，
求环管中心位置的磁场峰值。

\paragraph{习题 17.13}
图17.44所示环形螺线管沿高
\(h\)、内半径 \(R_1\)、外半径
\(R_2\) 的长方形截面的整个木环
均匀密绕 \(N\) 匝。
通入电流 \(I\)，忽略木材磁化与漏磁。
求环内及环外磁场分布，并求通过管截面的总磁通量。

\paragraph{习题 17.14}
两块平行的无限大金属薄板，
各有大小相等而方向相反的均匀面电流密度
\(j\)，求两板之间及两板之外的磁场分布。
\textit{这里 \(j\) 按原文“面电流密度”理解为
单位横向宽度的电流，单位 \(\mathrm{A/m}\)。}

\paragraph{习题 17.15}
无限长导体圆柱半径为 \(R\)，
通入沿轴向的总电流 \(I\)，
截面上各处电流密度均匀。
求柱内与柱外的磁场分布；
在长为 \(l\) 的一段圆柱内部，
求环绕中心轴线（经径向截面定义）的磁通量。

\paragraph{习题 17.16}
一根半径为 \(R\) 的长圆柱形导体，
在内部挖去与圆柱轴平行的半径为 \(r\)
的圆柱形空洞，两个轴线
\(O,O'\) 相距 \(d\)，截面见图17.45。
导体实体内部有均匀轴向电流密度
\(\boldsymbol J=J\hat{\boldsymbol z}\)。
求空洞内各处的磁场，证明其均匀。
\textit{空洞轴与原柱轴不是重合的。}

\paragraph{习题 17.17}
亥姆霍兹线圈由两根互相平行、
轴线共线的细圆线圈构成（图17.46）；
每线圈半径 \(R\)，两线圈轴向间距也为
\(R\)，线圈通同方向的等大电流 \(I\)。
（1）求共轴 \(z\) 线上任一点的磁感应强度；
（2）证明中点 \(z=0\) 处
\(dB/dz=0,\ d^2B/dz^2=0\)。

\paragraph{习题 17.18}
一个塑料圆盘半径为 \(R\)，
表面均匀分布电荷 \(q\)。
证明绕过盘心且垂直盘面的轴
以角速度 \(\omega\) 匀速转动时，
盘心磁场
\[
 \boxed{B=\frac{\mu_0\omega q}{2\pi R}}.
\]

\paragraph{习题 17.19}
平行板电容器的两板均为半径
\(5.0\,\mathrm{cm}\) 的圆导体片，
充电时两板间电场变化率为
\[
 \frac{dE}{dt}=1.0\times10^{12}\,
 \mathrm{\frac{V}{m\,s}}.
\]
求：（1）两板间的位移电流；
（2）板边缘处的磁感应强度
\(B\)。忽略边缘电场。

\paragraph{习题 17.20}
一对圆形板组成的平行板电容器，
电容 \(C=1.0\times10^{-12}\,\mathrm F\)，
加频率 \(f=50\,\mathrm{Hz}\)、
峰值电压 \(U_m=1.74\times10^5\,\mathrm V\)
的交变电压。
求极板之间位移电流的最大值。
